\section{Method: A Bio-Inspired Continuity Harness}
\label{sec:method}

\subsection{Motivation: Transport, Not Post-Hoc Moderation}

The primary vulnerability in tool-using language agents is the unauthorized transport of untrusted information across control boundaries. When an agent retrieves a web page, message, or API response, malicious instructions embedded in the payload can contaminate later plans and tool-call arguments without triggering immediate textual moderation. Safety therefore requires regulating pathway transport, not only classifying single inputs.

Our design is motivated by protein translocation. In biological systems, a peptide is folded into a transportable conformation, admitted through a selective channel, and stripped of transient signal sequences after crossing the membrane. This analogy yields three engineering principles. First, potentially adversarial content should be reorganized and stress-tested before it is trusted, because attacks often rely on carefully ordered fragments that appear rational only in their original sequence. Second, transport into a more consequential region should require directional authorization. Third, temporary routing markers should be removed downstream after they have served their transport role.

These principles lead to a three-stage harness---\textsc{Fold}, \textsc{Gate}, and \textsc{Cleave}---that forms a closed-loop transport protocol rather than a stack of independent filters. \textsc{Fold} organizes untrusted content, tests cross-step consistency, and emits a compact provenance tag. \textsc{Gate} authorizes the next transition according to direction and consequence: read-like continuations are treated as passive, while actions that affect external state require stronger authorization. \textsc{Cleave} performs the minimal repair authorized by Gate; a lightweight carry state then records what was removed and how the next step should be constrained. This separation preserves benign multi-step workflows while preventing injected instructions from becoming actions.

\begin{table}[t]
\centering
\small
\setlength{\tabcolsep}{3pt}
\caption{Compact view of the continuity harness.}
\label{tab:compact-design}
\begin{tabular}{p{0.18\columnwidth}p{0.26\columnwidth}p{0.45\columnwidth}}
\toprule
Stage & Analogue & Role \\
\midrule
\textsc{Fold} & chaperone & organize, reflect, and stress-test untrusted content before trust transfer \\
\textsc{Gate} & channel & authorize passive consumption or active externalization by continuity and sensitivity \\
\textsc{Cleave} & signal cut & remove transient transport markers and update lightweight carry state \\
\bottomrule
\end{tabular}
\end{table}

\subsection{Pipeline Overview}

The harness wraps an existing tool-using agent at two transport boundaries: when untrusted \emph{tool output} enters the agent loop, and when the agent proposes a \emph{tool call}. At each boundary the same protocol runs in order: Fold senses and tags the content, Gate decides whether the transition is passive, active, or blocked, Cleave performs authorized repair, and the carry state is updated. Fold consumes the user objective together with recent planning and review context; Gate and Cleave consume Fold's discrete provenance tag and the current carry level rather than recomputing Fold scores.

This ordering is deliberate. Injected content in a tool output is sensed and, when Gate permits, stripped before the model re-reads it. A later sensitive tool call is then judged under the repaired state rather than under stale pre-repair evidence. The deployed system, which we call \emph{BioHarness}, realizes this protocol with a learned provenance classifier, feature-attribution reordering, projected provenance features, inbound metadata repair, per-tool sensitivity filtering, and carry feedback. Continuity follows a Fold--Gate--Cleave orchestration profile that enables metadata stripping and goal-aligned continuation without reintroducing text-level reordering on live tool responses.

\subsection{Fold: Provenance Drift, Reflection, and Shuffle}

\textsc{Fold} structures untrusted input before validating semantic coherence. Static schemas alone are insufficient: an attacker can write a plan that looks well ordered while embedding a harmful action across fields or steps. Fold therefore combines canonical organization, shuffle-based coherence checking, reflection-style provenance sensing, and a compact provenance tagger.

\paragraph{Canonical organization.}
Typed intermediate content is parsed into fact, parameter, route, or metadata fragments. Malformed structures and missing required fields are rejected before trust transfer. The same parser is reused by Cleave when selecting removable chunks.

\paragraph{Shuffle coherence.}
Fold applies shuffle checks after canonical organization. We use two complementary variants depending on the evaluation setting.

\textit{Text reordering} (scripted regression suite) applies $K$ seeded within-type shuffles sequentially to the organized fragment sequence $S_t^{(0)}$,
\begin{equation}
    S_t^{(j)} = \mathrm{Shuffle}\!\left(S_t^{(j-1)};\,\pi_j\right),
    \qquad j=1,\ldots,K,
\end{equation}
and inspects the final reconstruction $S_t^{(K)}$. The default is $K=2$ (at most five under adaptive probing). Explicit rules flag nearby harmful and execution cues, contradictory claims, decomposed-and-reassembled instructions, suspicious schema fields, and gradual escalation. This suite verifies coverage of the listed constructions; it is not a general semantic invariance test.

\textit{Feature-attribution reordering} (live AgentDojo configuration) permutes text assigned to the review, tool-context, and proposed-action slots, reevaluates the provenance classifier, and measures the range of predicted scores. A change of at least $0.12$ triggers a warning when at least one reordered score is $0.18$ or higher. Separate path-consistency rules flag a proposed path or destination that conflicts with the user objective and prior review. The live configuration disables text reordering and enables this feature-space test.

\paragraph{Reflective Fold.}
Reflective Fold is not intended as a hard block that maximizes attack stop rate in isolation. Its role is to surface early geometric drift signals and feed them into continuity-aware authorization: when a proposed action remains goal-aligned, or when the next tool is a low-consequence read, the probe can soften lockdown rather than ratchet the episode into permanent refusal. Let $P(\cdot)$ map text into $d=64$ dimensions and $\mathrm{unit}(v)=v/\|v\|_2$. For user objective $u$, review context $r$, and proposed action slot $s$,
\begin{equation}
\begin{aligned}
q &= \mathrm{unit}(P(u)),\\
k_r &= \mathrm{unit}(P(r)), \qquad
k_s = \mathrm{unit}(P(s)),\\
\alpha_r &= q^\top k_r/\sqrt d,\qquad
\alpha_s = q^\top k_s/\sqrt d .
\end{aligned}
\end{equation}
The pairwise distance is written with a Householder reflection
\begin{equation}
    \rho_q(k)=k-2(q^\top k)q,\qquad
    \delta_{\mathrm{pair}}
    =\|\rho_q(k_s)-\rho_q(k_r)\|_2,
\end{equation}
which equals $\|k_s-k_r\|_2$ because $\rho_q$ is orthogonal. Operationally, therefore, reflection contributes a projected pairwise-drift feature rather than a new distance geometry. It is combined with $|\alpha_s-\alpha_r|$, sign changes across process segments, repetition without progress, and hesitation markers. The cross-step component is scored only when concrete values in the proposed action have token overlap below $0.08$ with the user objective. A composite score of $0.40$ raises a warning that Gate may soften or escalate depending on continuity and tool sensitivity, rather than an unconditional hard stop.

\paragraph{Provenance tagging.}
Geometric features, cross-step features, and hand-engineered provenance features enter a small two-layer classifier that emits a drift probability. At a tool boundary, a score of $0.20$ or higher recommends blocking. The $0.35$ threshold accumulates suspicious evidence produced before a tool boundary, or evidence labelled as untrusted carry; it is not a second ordered cutoff at the same boundary. Fold exports one of five handoff tags---aligned, injection watch, untrusted carry, feature-attribution drift, or projected provenance drift---which Gate consumes directly.

\subsection{Gate: Directional Authorization}

\textsc{Gate} distinguishes inbound consumption from outbound externalization. Reading or summarizing a repaired observation is passive. Sending a message, modifying shared state, making a reservation, or transferring value is active because it affects the environment.

Let $c_t\in\{0,1,2,3\}$ denote the carry level (none, watch, elevated, lockdown), with outbound sensitivity ceilings
\begin{equation}
    \Gamma(c_t)\in\{1.0,\,0.85,\,0.70,\,0.0\}.
\end{equation}
A fixed table assigns sensitivities to write-like tools. Let $z_t$ be Fold's discrete tag and $a_t$ the proposed action. The deployed policy can be summarized as
\begin{equation}
G_t =
\begin{cases}
\mathrm{passive}, & \text{inbound content is repairable},\\
\mathrm{block}, & z_t \text{ requires a hold and }
                  \mathrm{Sens}(a_t)>\Gamma(c_t),\\
\mathrm{active}, & \text{the outbound action may be reconsidered},\\
\mathrm{passive}, & \text{otherwise}.
\end{cases}
\end{equation}
Projected-provenance lockdown can impose a stricter hold than the generic sensitivity comparison. Conversely, a proposed action whose concrete argument values remain continuous with the user objective can soften stale projected-provenance evidence before Gate evaluates the action. Continuity is therefore a pre-authorization state update rather than a second score recomputed inside Gate.

\subsection{Cleave: Repair Without Over-Withholding}

\textsc{Cleave} removes transient transport material after Gate has authorized repair---injection envelopes, routing hints, and transport metadata---while preserving factual fields needed for the legitimate task. The implementation uses four repair scopes,
\begin{equation}
    \sigma_t\in
    \{\textsc{None},\textsc{Transit},\textsc{Metadata},\textsc{Strict}\},
\end{equation}
selected by Gate mode, carry level, and detected category. Passive inbound transport uses a narrow scope. A blocked transition uses \textsc{None}: held content is not repaired and silently passed. After repair, a release veto applies if the provenance score remains at least $0.45$. Successful stripping reduces stored contamination pressure without erasing the episode history that Gate still needs for consequential actions.

\subsection{Continuity and Carry State}
\label{sec:continuity}

Continuity asks whether the concrete values being externalized are supported by the user objective. The implementation extracts argument values such as recipients, message bodies, file names, and destinations while ignoring tool schema names and transport identifiers:
\begin{equation}
\mathrm{Cont}(u,a_t)
=
\max_{v\in\mathrm{Args}(a_t)}
\frac{|\mathrm{Tok}(u)\cap\mathrm{Tok}(v)|}
     {|\mathrm{Tok}(u)\cup\mathrm{Tok}(v)|}.
\end{equation}
The threshold is $\theta_c=0.08$.\footnote{Deployed BioHarness defaults used throughout our AgentDojo runs: continuity $\theta_c{=}0.08$; projected-provenance composite warning $0.40$; boundary classifier block $0.20$; pre-boundary carry $0.35$; projection dimension $d{=}64$; feature-reordering score range $0.12$ with maximum score $0.18$; carry ceilings $\Gamma{\in}\{1.0,0.85,0.70,0.0\}$; release veto at provenance score $0.45$; clear non-lockdown carry after three benign boundaries.} This narrow check does not grant broad permission merely because a recipient or entity appeared somewhere in the conversation.

Carry records whether an injection was observed and whether unresolved evidence remains. An injection envelope begins at watch level. Classifier drift raises the state to watch or elevated depending on the boundary; feature-attribution drift raises it to elevated when no injection is yet known; projected-provenance drift enters lockdown. Successful Cleave reduces the stored score, and three consecutive benign boundaries clear non-lockdown state. When a new stored episode begins, inherited lockdown is reduced to watch rather than applied indefinitely.

\section{Experimental Setup}

\subsection{Evaluation Settings}

We use two complementary settings.
\begin{enumerate}
    \item \textbf{AgentDojo interactive evaluation.} BioHarness is inserted at tool-output and tool-call boundaries in workspace, banking, and Slack environments. The fixed manifest contains 24 conditions per model: four attack framings, two injection tasks, and three environment suites. A banking condition contains six user-task episodes, a Slack condition two, and a workspace condition one. Component ablations reuse the same 24 DeepSeek conditions under full BioHarness and single-module removals.
    \item \textbf{Scripted mechanism checks.} A controlled replay suite examines repair and post-injection authorization. A held-out provenance audit measures in-distribution versus unseen attack constructions under the same BioHarness profile. A short text-reordering check verifies coverage of six designed constructions on the non-live text path.
\end{enumerate}

\subsection{Compared Defenses and Models}

Each AgentDojo condition evaluates the same task--attack pairs under no defense, Spotlighting with delimiters, a learned prompt-injection detector, and BioHarness. We evaluate DeepSeek-Chat (V3) and Llama-3-70B.

\subsection{Primary and Sensitivity Analyses}
\label{sec:stress-gate}

For each model, all 24 manifest conditions are retained for an unfiltered analysis. We additionally report a stress-conditioned analysis restricted to conditions where the attack succeeds against the undefended model:
\begin{equation}
    \mathrm{ASR}_{\mathrm{none}}\geq 10\%.
\end{equation}
Because $\mathrm{Security}=1-\mathrm{ASR}$, the equivalent condition $\mathrm{Security}_{\mathrm{none}}<90\%$ is not a separate criterion. The stress-conditioned set contains 18 DeepSeek and 16 Llama conditions. Selection depends only on the undefended baseline.

The primary tables use an equal-weight mean across experimental conditions so that six-episode banking conditions do not dominate one-episode workspace conditions. Because condition sizes differ, we also report an episode-weighted sensitivity analysis. Bootstrap intervals resample conditions and quantify variation across the tested manifest.

\subsection{Metrics}

\textbf{Security} is one minus attack success rate. \textbf{Utility} is the fraction of user-task episodes completed. Within each condition these are episode proportions; the reported macro result is the mean condition-level proportion. \textbf{Relative utility} is the ratio between a defense's aggregate utility and the no-defense aggregate on the same condition set.

The staged repair analysis replays one representative injected observation and scores five task-preservation checks plus one malicious action. Live component ablations report macro utility and security after removing one module from full BioHarness. The held-out provenance audit reports attack block rate and benign false-positive rate on scripted constructions.
